\section{Exercice de conception}

\begin{UPSTIexercice}{Alarme de dépassement de seuil}
    On simule la surveillance d'une grandeur mesurée périodiquement par un capteur (une nouvelle mesure est saisie à chaque tour de boucle).

    \textbf{Cahier des charges.} Le programme doit se comporter ainsi :
    \begin{itemize}
        \item Le programme demande à l'utilisateur de saisir successivement des mesures, une par une.
        \item La saisie s'arrête dès que l'utilisateur entre la valeur \texttt{-999} (valeur sentinelle, qui ne représente pas une mesure réelle).
        \item Si une mesure dépasse strictement 50, le programme affiche un message d'alerte.
        \item Lorsque la saisie est terminée, le programme affiche le nombre total de mesures qui ont déclenché une alerte.
    \end{itemize}

    \UPSTIquestion{Proposer un \textbf{pseudo-code} pour ce programme.}
\end{UPSTIexercice}

\begin{UPSTIprofOnlyEnv}
    \begin{UPSTIcorrectionP}{Alarme de dépassement de seuil}
        \begin{verbatim}
total_depassements <- 0
AFFICHER "Mesure (-999 pour arreter) : "
LIRE mesure

TANT QUE mesure != -999 FAIRE
    SI mesure > 50 ALORS
        AFFICHER "ALERTE : depassement de seuil"
        total_depassements <- total_depassements + 1
    FIN SI
    AFFICHER "Mesure (-999 pour arreter) : "
    LIRE mesure
FIN TANT QUE

AFFICHER "Nombre de depassements : ", total_depassements
        \end{verbatim}
    \end{UPSTIcorrectionP}
\end{UPSTIprofOnlyEnv}

\begin{UPSTIexercice}{Régulation d'un chauffage par hystérésis}
    On simule la régulation tout-ou-rien d'un chauffage d'appoint, asservi à une température mesurée périodiquement par un capteur (une nouvelle mesure est saisie à chaque tour de boucle).

    \textbf{Cahier des charges.} Le programme doit se comporter ainsi :
    \begin{itemize}
        \item Le programme demande à l'utilisateur de saisir successivement des températures mesurées (des nombres, en °C), une par une.
        \item La saisie s'arrête dès que l'utilisateur entre la valeur \texttt{-999} (valeur sentinelle, qui ne représente pas une température réelle).
        \item Le chauffage possède un état : \textbf{allumé} ou \textbf{éteint}. Avant la première mesure, il est éteint.
        \item Si la température mesurée est \textbf{strictement inférieure à 18}, le chauffage doit être allumé (s'il ne l'était pas déjà).
        \item Si la température mesurée est \textbf{strictement supérieure à 20}, le chauffage doit être éteint (s'il ne l'était pas déjà).
        \item Si la température mesurée est comprise entre 18 et 20 (bornes incluses), le chauffage \textbf{garde l'état qu'il avait} avant cette mesure : il ne doit pas changer d'état dans cette zone.
        \item Après chaque mesure, le programme affiche la température lue ainsi que le nouvel état du chauffage.
        \item Lorsque la saisie est terminée, le programme affiche le nombre total de fois où le chauffage a changé d'état (d'allumé à éteint, ou d'éteint à allumé) au cours de la simulation.
    \end{itemize}

    \UPSTIquestion{Proposer un \textbf{pseudo-code} pour ce programme.}
\end{UPSTIexercice}

\begin{UPSTIprofOnlyEnv}
    \begin{UPSTIcorrectionP}{Régulation d'un chauffage par hystérésis}
        La difficulté de conception ne vient pas de la boucle (arrêt sur sentinelle \texttt{-999}, comme dans l'exercice précédent), mais du fait que la décision « allumer ou éteindre » ne dépend pas \textbf{seulement} de la température qui vient d'être lue : elle dépend aussi de l'état \textbf{précédent} du chauffage, qu'il faut donc conserver d'une itération à l'autre. La zone intermédiaire \([18, 20]\) ne doit \textbf{rien changer} : un piège classique est d'écrire un \texttt{SI}/\texttt{SINON} à deux branches qui, sans le vouloir, force une décision (allumé ou éteint) même dans cette zone -- il faut bien une structure à trois cas, dont un qui ne fait rien à l'état.

        \begin{verbatim}
chauffage_allume <- FAUX
nombre_changements <- 0

AFFICHER "Temperature (-999 pour arreter) : "
LIRE temperature

TANT QUE temperature != -999 FAIRE
    SI temperature < 18 ALORS
        nouvel_etat <- VRAI
    SINON SI temperature > 20 ALORS
        nouvel_etat <- FAUX
    SINON
        nouvel_etat <- chauffage_allume  % zone morte : aucun changement
    FIN SI

    SI nouvel_etat != chauffage_allume ALORS
        nombre_changements <- nombre_changements + 1
    FIN SI
    chauffage_allume <- nouvel_etat

    AFFICHER "Temperature : ", temperature, " -> chauffage ", chauffage_allume
    AFFICHER "Temperature (-999 pour arreter) : "
    LIRE temperature
FIN TANT QUE

AFFICHER "Nombre de changements d'etat : ", nombre_changements
        \end{verbatim}

        \textbf{Pourquoi deux seuils plutôt qu'un seul ?} Avec un seuil unique (par exemple : allumer si $< 19$, éteindre si $> 19$), une température qui oscille tout près de 19°C ferait changer d'état le chauffage à chaque mesure -- ce qu'on appelle le \textbf{pompage}. C'est pour l'éviter que les systèmes de régulation tout-ou-rien utilisent une \textbf{bande morte} (ici, l'intervalle \([18, 20]\)) entre le seuil de mise en marche et le seuil d'arrêt : c'est le principe de l'\textbf{hystérésis}, que vous retrouverez en automatisme.
    \end{UPSTIcorrectionP}
\end{UPSTIprofOnlyEnv}
